% Original analytic teaching diagram, 2026-11-12.
% Where the parameter enters J(theta) = E[f_theta(X)].  Each column: source of randomness (top), sample X,
% function value f_theta(X), expectation J(theta).  Red theta nodes mark every place the parameter enters.
% (a) fixed distribution P: theta only in f.  (b) generation map X = T_theta(eps) with fixed noise law Q:
% theta in T and in f.  (c) distribution P_theta (probability weights): theta in P_theta and in f.
\definecolor{VB}{HTML}{7188A3}\definecolor{VBF}{HTML}{E5EAF0}
\definecolor{VG}{HTML}{829C8B}\definecolor{VGF}{HTML}{E8EEE9}
\definecolor{VY}{HTML}{B4A16F}\definecolor{VYF}{HTML}{F1ECDD}
\definecolor{VR}{HTML}{A7787D}\definecolor{VRF}{HTML}{F0E5E6}
\definecolor{VP}{HTML}{9588A5}\definecolor{VPF}{HTML}{ECE8F0}
\definecolor{VI}{HTML}{3B4252}\definecolor{VM}{HTML}{6B7280}\definecolor{VL}{HTML}{4C4D4F}
\begin{tikzpicture}[x=1mm,y=1mm,>=Stealth,font=\figurefont\fontsize{8}{10}\selectfont,
lab/.style={text=VM,font=\figurefont\fontsize{7.5}{9}\selectfont},
dist/.style={rectangle,rounded corners=1mm,draw=VB,line width=.7pt,fill=VBF,minimum width=12mm,minimum height=5.5mm,inner sep=1.2pt,text=VI},
box/.style={rectangle,draw=VL,line width=.5pt,fill=white,minimum width=12mm,minimum height=5.5mm,inner sep=1.2pt,text=VI},
th/.style={circle,fill=VRF,draw=VR,line width=.6pt,inner sep=.6pt,minimum size=4.6mm,text=VI},
ar/.style={->,VL,line width=.55pt,shorten >=.4mm,shorten <=.4mm},
ra/.style={->,VR,line width=.75pt,shorten >=.4mm,shorten <=.4mm}]
% ===== (a) =====
\node[dist] (a0) at (0,0) {$P$};
\node[box] (a1) at (0,-12) {$X$};
\node[box] (a2) at (0,-24) {$f_{\bm{\theta}}(X)$};
\node[box] (a3) at (0,-36) {$J(\bm{\theta})$};
\node[th] (at) at (17,-24) {$\bm{\theta}$};
\draw[ar] (a0)--(a1) node[lab,midway,left=.8mm] {抽样};
\draw[ar] (a1)--(a2);
\draw[ar] (a2)--(a3) node[lab,midway,left=.8mm] {平均};
\draw[ra] (at)--(a2);
\node[lab,text=VR,anchor=south] at (11.5,-23) {函数值};
\node[text=VI] at (4,-43) {(a)};
% ===== (b) =====
\begin{scope}[shift={(38,0)}]
\node[dist] (b0) at (0,0) {$Q$};
\node[box] (b1) at (0,-12) {$X=T_{\bm{\theta}}(\varepsilon)$};
\node[box] (b2) at (0,-24) {$f_{\bm{\theta}}(X)$};
\node[box] (b3) at (0,-36) {$J(\bm{\theta})$};
\node[th] (bt) at (21,-18) {$\bm{\theta}$};
\draw[ar] (b0)--(b1) node[lab,midway,left=.8mm] {抽 $\varepsilon$};
\draw[ar] (b1)--(b2);
\draw[ar] (b2)--(b3);
\draw[ra] (bt)--(b1.east);
\draw[ra] (bt)--(b2.east);
\node[lab,text=VR,anchor=south west] at (13,-9) {生成映射};
\node[lab,text=VR,anchor=north west] at (13,-25.5) {函数值};
\node[text=VI] at (4,-43) {(b)};
\end{scope}
% ===== (c) =====
\begin{scope}[shift={(78,0)}]
\node[dist] (c0) at (0,0) {$P_{\bm{\theta}}$};
\node[box] (c1) at (0,-12) {$X$};
\node[box] (c2) at (0,-24) {$f_{\bm{\theta}}(X)$};
\node[box] (c3) at (0,-36) {$J(\bm{\theta})$};
\node[th] (ct) at (16,-12) {$\bm{\theta}$};
\draw[ar] (c0)--(c1);
\draw[ar] (c1)--(c2);
\draw[ar] (c2)--(c3);
\draw[ra] (ct)--(c0.east);
\draw[ra] (ct)--(c2.east);
\node[lab,text=VR,anchor=south west] at (10,-3) {概率权重};
\node[lab,text=VR,anchor=north west] at (10.5,-20.5) {函数值};
\node[text=VI] at (4,-43) {(c)};
\end{scope}
\end{tikzpicture}
